\section{Related Work}
\label{sec:related}

Our work lies at the intersection of three research threads: KG embedding evaluation, uncertainty quantification in knowledge graphs, and multiple hypothesis testing in machine learning.

\subsection{KG Embedding and Link Prediction Evaluation}

Knowledge graph embedding models map entities and relations into continuous vector spaces and score candidate triples via a learned scoring function~\cite{wang2017knowledge}. Large-scale benchmarking efforts~\cite{ali2021pykeen,rossi2021knowledge} have standardized dataset splits and hyperparameter tuning, and diverse model families---translation-based, tensor factorization, convolutional, GNN-based, hyperbolic, and LLM-enhanced---have been developed~\cite{bordes2013translating,sun2019rotate,trouillon2016complex,dettmers2018convolutional,balazevic2019tucker,vashishth2020composition,zhu2021neural,le2023knowledge,lu2024deep,fang2025knowledge,gao2025unifying}. The standard evaluation protocol, established by~\cite{bordes2013translating}, measures MRR and Hits@$K$ on held-out test triples under the filtered setting, and has enabled reproducible comparison across this diverse model landscape.

\textbf{Limitations of ranking-based evaluation.}
A growing body of work has identified critical shortcomings of aggregate ranking metrics.
\cite{reimers2017reporting} argued that score distributions reveal substantially more about model behavior than point estimates, yet did not propose a statistical inference framework.
\cite{safavi2020evaluating} and~\cite{wang2023evaluation} showed that KG completion benchmarks contain systematic biases---including entity degree bias~\cite{shomer2023toward} and hyperparameter sensitivity~\cite{lloyd2023assessing}---that inflate apparent performance, while~\cite{liu2022understanding} systematically analyzed fundamental limitations of standard evaluation protocols without proposing a remedial statistical framework.
These critiques converge on a common insight---a single scalar cannot capture the full picture of model quality---but stop short of the question we answer: \emph{which} individual predictions are reliable and \emph{what fraction} of the test set a model handles at above-chance level?

\textbf{Calibration and confidence estimation.}
In parallel, several authors have noted that KG embedding scores are not calibrated as probabilities and cannot be directly interpreted as confidence measures.
\cite{tabacof2020probability} demonstrated this calibration failure and proposed Platt scaling as a post-hoc correction, but calibration alone does not provide formal error rate control.
\cite{jarnac2025uncertainty} surveyed the broader landscape of uncertainty management in KG construction and identified the lack of calibrated confidence estimates as a critical open gap.
Our work directly addresses this gap: rather than calibrating scores to probabilities, we construct valid $p$-values via permutation-based negative sampling, enabling formal FDR control procedures to provide per-prediction reliability guarantees with controlled error rates.

\subsection{Uncertainty Quantification in Knowledge Graphs}

Several lines of work have pursued uncertainty quantification for KG reasoning.

\textbf{Gaussian and Bayesian approaches.}
\cite{he2015learning} proposed KG2E, which represents entities and relations as multivariate Gaussian distributions, enabling variance-based uncertainty estimation at the entity and relation level.
\cite{wu2025uncertain} introduced semi-supervised confidence distribution learning for uncertain KG completion.
These methods produce uncertainty estimates during training but provide no formal statistical guarantee on the error rate of downstream predictions---a limitation we overcome by constructing valid $p$-values that directly control the false discovery rate.

\textbf{Conformal prediction for KGs.}
A recent and closely related line of work applies conformal prediction to KG link prediction.
\cite{marandon2024conformal} combined conformal prediction with link prediction to achieve FDR control, published in a statistics journal---this is the closest prior work to ours.
\cite{zhu2025conformal} applied conformal prediction to KG embeddings specifically, producing prediction sets with finite-sample coverage guarantees.
\cite{zhu2025certainty} (UnKGCP) extended this direction by introducing a tailored nonconformity measure for sharper prediction intervals.
\cite{ni2025towards} coupled conformal prediction with LLM-based KG reasoning for error rate control.

Conformal prediction and the FDR framework address complementary but distinct questions:

\begin{center}
\begin{tabular}{p{0.44\textwidth} p{0.44\textwidth}}
\toprule
\textbf{Conformal Prediction (Prior Work)} & \textbf{FDR Framework (This Work)} \\
\midrule
Produces prediction sets with coverage guarantees & Produces per-triple significance decisions with error rate control \\
Single guarantee level (e.g., $1 - \alpha$ coverage) & Multi-level diagnostics (global $\hat{\pi}_0$, per-relation FDR, per-triple local FDR) \\
No direct estimate of noise proportion & $\hat{\pi}_0$ estimates fraction of unpredictable triples \\
No per-relation stratification & Relation-stratified FDR reveals heterogeneity invisible to aggregate metrics \\
\bottomrule
\end{tabular}
\end{center}

Conformal methods answer ``what is the smallest prediction set with guaranteed coverage?'' while FDR answers ``how many individual predictions are statistically significant, which relations are reliable, and what fraction of the test set is noise?''
Both perspectives are valuable; we provide the first systematic FDR-based evaluation protocol for KG completion.

\textbf{Evaluation reform and predictive multiplicity.}
Broader evaluation concerns in KG completion reinforce the need for statistical approaches.
\cite{zhu2024predictive} demonstrated that KG embedding models exhibit predictive multiplicity: re-training with different random seeds produces substantially different rankings, directly motivating statistical methods that quantify ranking variability rather than relying on single-run point estimates.
Comprehensive surveys~\cite{shen2022comprehensive,li2023knowledge,jarnac2025uncertainty} have consistently called for more rigorous evaluation methodologies, and recent re-evaluation studies~\cite{sun2020reevaluation,schwertmann2023modelagnostic} have shown that standard metrics can conceal systematic deficiencies.

\subsection{Multiple Hypothesis Testing in Machine Learning}

The FDR framework~\cite{benjamini1995controlling} and its extensions~\cite{storey2002direct,storey2003statistical,efron2004large,efron2007size} have been foundational in genomics and have recently gained traction in machine learning.
\cite{demsar2006statistical} introduced statistical comparisons of classifiers over multiple datasets to the ML community, advocating for Friedman tests and Nemenyi post-hoc analysis, though this work focused on aggregate classifier comparison rather than per-instance reliability assessment.
Selective inference~\cite{taylor2015statistical,lee2016exact} addresses the ``data snooping'' problem by providing valid post-selection $p$-values in adaptive data analysis, operating in a fundamentally different setting from the batch evaluation problem we address.
The knockoff filter~\cite{barber2015controlling,candes2018panning} offers finite-sample FDR control for variable selection under controlled dependence structures; its adaptation to graph-structured KG data is an open direction.
More recent advances include graph-aware multiple testing and conformalized FDR approaches~\cite{angelopoulos2023conformal,angelopoulos2024prediction,bates2024cross,marandon2024fdr}, which share our goal of rigorous error control but differ in the inferential target---prediction sets versus per-triple significance decisions.

Despite this rich statistical toolkit, its systematic application to KG completion evaluation remains largely unexplored.
Our work bridges this gap: we adapt empirical $p$-value construction, BH FDR control~\cite{benjamini1995controlling,benjamini2001control}, Storey's $q$-value estimation, and Efron's empirical Bayes local FDR to the specific structure of KG link prediction, producing a three-level evaluation protocol that complements existing ranking-based metrics.
